% ═══════════════════════════════════════════════════════════════
% ML-01b: Newton 3 + Lageenergie – MUSTERLÖSUNG
% Physik Klasse 10 MR
% ═══════════════════════════════════════════════════════════════

\documentclass[11pt, a4paper]{article}

\usepackage[utf8]{inputenc}
\usepackage[T1]{fontenc}
\usepackage[ngerman]{babel}

\usepackage{helvet}
\renewcommand{\familydefault}{\sfdefault}

\usepackage{microtype}
\usepackage[margin=1.5cm, top=2cm, bottom=1.5cm]{geometry}
\usepackage{enumitem}
\usepackage{tabularx}
\usepackage{booktabs}
\usepackage{array}
\usepackage{tikz}
\usetikzlibrary{calc, arrows.meta}

\usepackage{amsmath, amssymb}
\usepackage{siunitx}
\sisetup{locale = DE, per-mode = symbol}

\usepackage{fancyhdr}
\usepackage{lastpage}
\usepackage{needspace}

\usepackage{xcolor}
\usepackage{tcolorbox}
\tcbuselibrary{skins}

% ═══════════════════════════════════════════════════════════════
% FARBEN
% ═══════════════════════════════════════════════════════════════

\definecolor{physik}{HTML}{2c5aa0}
\definecolor{hellgrau}{HTML}{F5F5F5}
\definecolor{mittelgrau}{HTML}{888888}
\definecolor{rahmen}{HTML}{CCCCCC}
\definecolor{kraftrot}{HTML}{c62828}
\definecolor{kraftblau}{HTML}{1565c0}
\definecolor{loesung}{HTML}{c62828}

\colorlet{fachfarbe}{physik}

% ═══════════════════════════════════════════════════════════════
% VARIABLEN
% ═══════════════════════════════════════════════════════════════

\newcommand{\fach}{Physik}
\newcommand{\thema}{Newton 3 + Lageenergie}
\newcommand{\klasse}{10 MR}
\newcommand{\stunde}{01b}
\newcommand{\maxpunkte}{20}

% ═══════════════════════════════════════════════════════════════
% KOPF- UND FUSSZEILE
% ═══════════════════════════════════════════════════════════════

\pagestyle{fancy}
\fancyhf{}
\renewcommand{\headrulewidth}{0.4pt}
\renewcommand{\footrulewidth}{0pt}
\lhead{\small \fach{} Klasse \klasse{} | \thema{} | \textcolor{loesung}{\textbf{MUSTERLÖSUNG}}}
\rhead{\small Name: \underline{\hspace{3.5cm}}}
\cfoot{\small\textcolor{mittelgrau}{Seite \thepage{} von \pageref{LastPage}}}

% ═══════════════════════════════════════════════════════════════
% BOXEN
% ═══════════════════════════════════════════════════════════════

\newtcolorbox{merkkasten}[1][]{
    colback=hellgrau,
    colframe=hellgrau,
    coltitle=black,
    colbacktitle=hellgrau,
    boxrule=0pt,
    arc=0pt,
    left=8pt, right=8pt, top=8pt, bottom=6pt,
    borderline north={2pt}{0pt}{fachfarbe},
    before skip=6pt, after skip=6pt,
    fonttitle=\bfseries,
    #1
}

\newtcolorbox{formelkasten}{
    colback=white,
    colframe=fachfarbe,
    boxrule=1.5pt,
    arc=0pt,
    left=8pt, right=8pt, top=6pt, bottom=6pt,
    before skip=4pt, after skip=4pt
}

% ═══════════════════════════════════════════════════════════════
% BEFEHLE
% ═══════════════════════════════════════════════════════════════

\newcommand{\luecke}[1]{\underline{\hspace{#1}}}
\newcommand{\lsg}[1]{\textcolor{loesung}{\textbf{#1}}}
\newcommand{\punktebox}[1]{\hfill\small\textcolor{mittelgrau}{[\textcolor{loesung}{#1}/#1]}}

\newcommand{\aufgabe}[1]{%
    \needspace{4\baselineskip}%
    \vspace{10pt}%
    \noindent\textbf{\textcolor{fachfarbe}{Aufgabe #1}}%
    \vspace{4pt}%
    \nopagebreak[4]%
}

\newcommand{\operator}[1]{\textbf{\textcolor{fachfarbe}{#1}}}

\newlist{teil}{enumerate}{2}
\setlist[teil,1]{label=\alph*), leftmargin=1.8em, itemsep=3pt, parsep=0pt, topsep=3pt}

\newcommand{\antwortzeilen}[1]{%
    \par\vspace{4pt}%
    \foreach \n in {1,...,#1}{%
        \noindent\rule{\linewidth}{0.3pt}\vspace{6pt}%
    }%
}

\setlength{\parindent}{0pt}
\setlength{\parskip}{4pt}

% ═══════════════════════════════════════════════════════════════
% DOKUMENT
% ═══════════════════════════════════════════════════════════════

\begin{document}

% ═══════════════════════════════════════════════════════════════
% HEADER
% ═══════════════════════════════════════════════════════════════

\begin{center}
    {\Large\bfseries\textcolor{fachfarbe}{\thema}}\\[0.2em]
    {\small Arbeitsblatt | \fach{} Klasse \klasse{} | Stunde \stunde{} | \textcolor{loesung}{\textbf{MUSTERLÖSUNG}}}
\end{center}
\vspace{-0.3em}
\hrule
\vspace{0.6em}

% ═══════════════════════════════════════════════════════════════
% TEIL 1: NEWTON 3
% ═══════════════════════════════════════════════════════════════

\textbf{\large Teil 1: Wechselwirkungsgesetz (Newton 3)}

\begin{merkkasten}[title={\textbf{Kasten 1: Wechselwirkungsgesetz}}]
\textbf{Actio = Reactio}

Kräfte treten immer \lsg{paarweise} auf.

\vspace{6pt}
\textbf{Drei Eigenschaften jedes Kräftepaares:}
\begin{enumerate}[leftmargin=2em, itemsep=2pt]
    \item Gleich \lsg{groß} (|$F_1$| = |$F_2$|)
    \item Entgegengesetzt \lsg{gerichtet}
    \item Wirken an \lsg{verschiedenen} Körpern!
\end{enumerate}
\end{merkkasten}

\aufgabe{1 – Kräftepaare zeichnen}

\operator{Zeichne} in jede Skizze die beiden Kräfte des Kräftepaares als Pfeile ein. Beschrifte die Pfeile!

\begin{teil}
\item \textbf{Schwimmer} ($\star$) \punktebox{2}

\begin{center}
\begin{tikzpicture}[scale=0.9]
    % Wasser
    \fill[blue!10] (0,0) rectangle (8,1.8);
    \node[blue!60] at (7.5, 0.3) {\tiny Wasser};
    % Strichmännchen (schwimmend, horizontal)
    % Kopf
    \draw[thick] (5.5,1.2) circle (0.3);
    % Körper
    \draw[thick] (5.2,1.2) -- (3.5,1.2);
    % Arm vorne (stößt Wasser nach hinten)
    \draw[thick] (4.2,1.2) -- (3.2,0.6);
    \draw[thick] (3.2,0.6) -- (2.8,0.5); % Hand
    % Arm hinten
    \draw[thick] (4.2,1.2) -- (4.8,1.7);
    % Beine
    \draw[thick] (3.5,1.2) -- (3.0,0.8);
    \draw[thick] (3.5,1.2) -- (3.0,1.5);
    % LÖSUNG: Kräftepfeile
    \draw[-{Stealth[length=3mm]}, line width=1.5pt, loesung] (2.8,0.5) -- (1.3,0.5);
    \node[loesung, above] at (2.0, 0.6) {\small $F_\text{Hand→Wasser}$};
    \draw[-{Stealth[length=3mm]}, line width=1.5pt, loesung] (3.2,0.6) -- (4.7,1.0);
    \node[loesung, above] at (4.2, 1.1) {\small $F_\text{Wasser→Hand}$};
\end{tikzpicture}
\end{center}
\vspace{0.3em}

\item \textbf{Rakete} ($\star\star$) \punktebox{2}

\begin{center}
\begin{tikzpicture}[scale=0.8]
    % Rakete - Nase zeigt nach OBEN
    % Rumpf
    \draw[thick, fill=gray!60] (2.5,0.5) rectangle (3.5,2.5);
    % Spitze (Nase oben)
    \draw[thick, fill=gray!50] (2.5,2.5) -- (3,3.5) -- (3.5,2.5) -- cycle;
    % Düse/Fins unten
    \draw[thick, fill=red!60] (2.3,0.5) -- (2.5,0.5) -- (2.5,0) -- cycle;
    \draw[thick, fill=red!60] (3.7,0.5) -- (3.5,0.5) -- (3.5,0) -- cycle;
    % Flamme/Gas
    \draw[thick, fill=orange!50] (3,-0.5) ellipse (0.4 and 0.25);
    \draw[thick, fill=yellow!70] (3,-0.4) ellipse (0.2 and 0.15);
    \node at (3,-1) {\tiny Gas};
    % LÖSUNG: Kräftepfeile
    \draw[-{Stealth[length=3mm]}, line width=1.5pt, loesung] (3,-0.5) -- (3,-1.8);
    \node[loesung, right] at (3.1, -1.2) {\small $F_\text{Rakete→Gas}$};
    \draw[-{Stealth[length=3mm]}, line width=1.5pt, loesung] (3,2.5) -- (3,4.0);
    \node[loesung, right] at (3.1, 3.3) {\small $F_\text{Gas→Rakete}$};
\end{tikzpicture}
\end{center}
\vspace{0.3em}

\item \textbf{Auto auf Straße} ($\star\star\star$) \punktebox{2}

\begin{center}
\begin{tikzpicture}[scale=0.8]
    % Straße
    \fill[gray!40] (0,-0.3) rectangle (8,0);
    % Auto
    \draw[thick, fill=blue!30] (2,0.3) rectangle (6,1.2);
    \draw[thick, fill=blue!30] (2.5,1.2) rectangle (5.5,1.8);
    % Räder
    \draw[thick, fill=gray!70] (2.8,0.15) circle (0.35);
    \draw[thick, fill=gray!70] (5.2,0.15) circle (0.35);
    % LÖSUNG: Kräftepfeile
    \draw[-{Stealth[length=3mm]}, line width=1.5pt, loesung] (2.8,-0.2) -- (1.5,-0.2);
    \node[loesung, below] at (2.15, -0.3) {\small $F_\text{Reifen→Straße}$};
    \draw[-{Stealth[length=3mm]}, line width=1.5pt, loesung] (5.2,0.5) -- (6.5,0.5);
    \node[loesung, above] at (5.85, 0.6) {\small $F_\text{Straße→Reifen}$};
\end{tikzpicture}
\end{center}

\end{teil}

\vspace{0.8em}

% ═══════════════════════════════════════════════════════════════
% TEIL 2: LAGEENERGIE
% ═══════════════════════════════════════════════════════════════

\textbf{\large Teil 2: Lageenergie}

\begin{merkkasten}[title={\textbf{Kasten 2: Lageenergie (potenzielle Energie)}}]
\begin{formelkasten}
\centering
\Large $E_\text{pot} = m \cdot g \cdot h$
\end{formelkasten}

\vspace{6pt}
\begin{tabular}{@{}ll@{}}
$E_\text{pot}$ & Lageenergie in \lsg{Joule (J)} \\[4pt]
$m$ & Masse in \lsg{kg} \\[4pt]
$g$ & Ortsfaktor: \textbf{\lsg{10} m/s²} (auf der Erde) \\[4pt]
$h$ & Höhe in \lsg{m (Meter)}
\end{tabular}
\end{merkkasten}

\begin{merkkasten}[title={\textbf{Kasten 3: Formel-Dreieck}}]
\begin{center}
\begin{tikzpicture}[scale=1]
    % Dreieck-Box
    \draw[thick, fachfarbe] (0,0) rectangle (3,1);
    \draw[thick, fachfarbe] (0,1) rectangle (3,2);
    \draw[thick, fachfarbe] (1.5,0) -- (1.5,1);
    % Beschriftungen
    \node at (1.5,1.5) {\large $E_\text{pot}$};
    \node at (0.75,0.5) {\large $m$};
    \node at (2.25,0.5) {\large $g \cdot h$};
\end{tikzpicture}
\end{center}

\vspace{4pt}
\textbf{Umstellen:}
\begin{itemize}[leftmargin=2em, itemsep=2pt]
    \item $E_\text{pot} = m \cdot g \cdot h$
    \item $m = E_\text{pot} \,/\, (g \cdot h)$
    \item $h = E_\text{pot} \,/\, (m \cdot g)$
\end{itemize}
\end{merkkasten}

\newpage

\aufgabe{2 – Lageenergie berechnen}

\operator{Berechne} die Lageenergie. Verwende $g = 10$ m/s².

\begin{teil}
\item ($\star$) Ein Ball ($m = 2$ kg) liegt auf einem Tisch ($h = 1$ m). \punktebox{2}

\vspace{0.3em}
$E_\text{pot} = $ \lsg{$m \cdot g \cdot h = 2\,\text{kg} \cdot 10\,\text{m/s}^2 \cdot 1\,\text{m} = 20\,\text{J}$}

\vspace{0.8em}

\item ($\star\star$) Ein Wassereimer ($m = 10$ kg) steht auf einem Dach ($h = 5$ m). \punktebox{2}

\vspace{0.3em}
$E_\text{pot} = $ \lsg{$m \cdot g \cdot h = 10\,\text{kg} \cdot 10\,\text{m/s}^2 \cdot 5\,\text{m} = 500\,\text{J}$}

\vspace{0.8em}

\item ($\star\star$) Ein Körper hat $E_\text{pot} = 400$ J in $h = 5$ m Höhe. \operator{Berechne} die Masse. \punktebox{2}

\vspace{0.3em}
$m = $ \lsg{$E_\text{pot} / (g \cdot h) = 400\,\text{J} / (10\,\text{m/s}^2 \cdot 5\,\text{m}) = 400 / 50 = 8\,\text{kg}$}

\end{teil}

\vspace{0.8em}

\aufgabe{3 – Formel umstellen}

($\star\star$) Ein 20-kg-Körper hat eine Lageenergie von 1000 J.

\operator{Berechne}, in welcher Höhe sich der Körper befindet. \punktebox{4}

\vspace{0.3em}
\textit{Gegeben:} \lsg{$m = 20\,\text{kg}$, $E_\text{pot} = 1000\,\text{J}$, $g = 10\,\text{m/s}^2$}

\vspace{0.3em}
\textit{Gesucht:} \lsg{$h$}

\vspace{0.3em}
\textit{Lösung:}

\lsg{$h = E_\text{pot} / (m \cdot g) = 1000\,\text{J} / (20\,\text{kg} \cdot 10\,\text{m/s}^2) = 1000 / 200 = 5$}

\vspace{0.3em}
\textit{Antwort:} $h = $ \lsg{5 m}

\vspace{1em}

\aufgabe{4 – Transfer: Der Staudamm}

($\star\star\star$) \operator{Erkläre}, warum ein Staudamm Energie speichern kann. \punktebox{4}

Verwende die Begriffe: \textit{Masse, Höhe, Lageenergie, $E_\text{pot} = m \cdot g \cdot h$}

\vspace{0.5em}
\textcolor{loesung}{\textbf{Musterlösung:}}

\textcolor{loesung}{Ein Staudamm speichert Energie, weil das Wasser \textbf{Lageenergie} besitzt.}

\textcolor{loesung}{Das Wasser hat eine große \textbf{Masse} $m$ und befindet sich auf einer bestimmten \textbf{Höhe} $h$ oberhalb des Tals.}

\textcolor{loesung}{Nach der Formel $E_\text{pot} = m \cdot g \cdot h$ steigt die Lageenergie mit der Masse und der Höhe.}

\textcolor{loesung}{Wenn das Wasser abgelassen wird, fällt es hinab und die Lageenergie wird in Bewegungsenergie umgewandelt, die dann eine Turbine antreibt.}

\textcolor{loesung}{Je mehr Wasser (größere Masse) und je höher der Staudamm (größere Höhe), desto mehr Energie kann gespeichert werden.}

\vspace{1em}

% ═══════════════════════════════════════════════════════════════
% FOOTER
% ═══════════════════════════════════════════════════════════════

\vfill
\hrule
\vspace{0.3em}
\begin{center}
\small\textcolor{mittelgrau}{Gesamt: \textcolor{loesung}{\textbf{20}} / \maxpunkte{} BE}
\end{center}

\end{document}
